% Fichier complété par tools/complete_missing_corrections.py.
% Source : CIN/CIN-02-VitesseAcceleration/16_Poussoir/16_Poussoir.tex
% Statut : rédigé (3 question(s)).
\begin{corrigebox}[Corrigé rédigé]
\CorrigeQuestion{1}
Le graphe comprend le bâti 0, la came 1 en pivot d'axe $(A,\vec k_0)$
                avec 0, le poussoir 2 en glissière suivant l'axe vertical avec 0, et un
                contact came--galet entre 1 et 2.
\CorrigeQuestion{2}
Pour $\theta=\pi/4$, on fait tourner la came de $45^\circ$ dans le
                sens positif puis on translate le poussoir jusqu'à retrouver la tangence
                avec le profil. La normale commune passe par le centre du galet.
\CorrigeQuestion{3}
Pour $\theta=-\pi/4$, la même construction est effectuée dans le sens
                opposé. La course obtenue est la valeur de la loi $\lambda=f(-\pi/4)$.
\end{corrigebox}
